%% ma: L10-B04-0001
%% lop: 10
%% chuong: 2
%% bai: 4
%% dang: 10.4.1
%% loai: TLN
%% muc: VD
%% dap_an: "5,8"
%% nguon: "(NBV)-ĐỀ SỐ 7-MỖI NGÀY 1 ĐỀ THI 2025.pdf (D07, MỖI NGÀY 1 ĐỀ - Đề số 7), Phần III Câu 1"
%% kien_thuc: "Bài 4 (hệ bất phương trình bậc nhất hai ẩn, miền nghiệm, giá trị lớn nhất của biểu thức)"
%% trang_thai: cho-duyet
%% ly_do: "Dạng toán mới đề xuất 10.4.1 (Quy hoạch tuyến tính: tối đa lợi nhuận đóng hai loại thuyền); chờ giáo viên duyệt dạng."
%% ghi_chu: "Bỏ ảnh minh họa xưởng đóng thuyền (không có dữ kiện toán). Đáp án gốc 5,8 đúng (kiểm bằng sympy/duyệt đỉnh). Lưu ý bảng đáp án gốc hoán đổi câu 5 và câu 6."
\begin{ex}
Một cơ sở đóng thuyền thủ công cần $10$ giờ lao động để đóng một thuyền loại A và $15$ giờ lao động để đóng một thuyền loại B. Mỗi tuần cơ sở bố trí được tối đa $120$ giờ lao động cho việc đóng hai loại thuyền này. Qua thực tế, người ta thấy mỗi tuần cơ sở bán được tối đa $6$ thuyền loại A và tối thiểu $2$ thuyền loại B. Mỗi thuyền loại A, loại B cho lợi nhuận lần lượt là $0{,}5$ triệu đồng và $0{,}7$ triệu đồng. Qua thời gian tìm hiểu, chủ cơ sở đã tìm được số lượng nên đóng thuyền mỗi loại để có thể thu được lợi nhuận cao nhất là $T$ (triệu đồng). Tìm $T$?
\shortans{5,8}
\loigiai{Gọi $x$, $y$ lần lượt là số thuyền loại A, loại B đóng trong một tuần. Ta có hệ $0\le x\le6$, $y\ge2$, $10x+15y\le120\Leftrightarrow2x+3y\le24$. Miền nghiệm là tứ giác có các đỉnh $(0;2)$, $(6;2)$, $(6;4)$, $(0;8)$. Lợi nhuận $T=0{,}5x+0{,}7y$ tại các đỉnh lần lượt là $1{,}4$; $4{,}4$; $5{,}8$; $5{,}6$. Giá trị lớn nhất là $5{,}8$ tại $(6;4)$ (là cặp số nguyên hợp lệ).}
\end{ex}
